% Ad Familiares 1.4 (ad-familiares-01-04)
% Source: https://raw.githubusercontent.com/PerseusDL/canonical-latinLit/master/data/phi0474/phi056/phi0474.phi056.perseus-lat1.xml
% Fetched by scripts/fetch_latin.py

% Dateline (Perseus): Scr. Romae circiter xv Kal. Febr. a. 698 (56).

\ciceroLetterOpener{M. CICERO S.D. P. LENTVLO PROCOS.}

\ciceroSection{1} A. d. xvi K. Febr. cum in senatu pulcherrime staremus, quod iam illam sententiam Bibuli de tribus legatis pridie eius diei fregeramus, unumque certamen esset relictum sententia Volcaci, res ab adversariis nostris extracta est variis calumniis; causam enim frequenti senatu non magna varietate magnaque invidia eorum, qui a te causam regiam alio traferebant, obtinebamus. eo die acerbum habuimus Curionem, Bibulum multo iustiorem, paene etiam amicum; Caninius et Cato negarunt se legem ullam ante comitia esse laturos. senatus haberi ante K. Februarias per legem Pupiam, id quod scis, non potest neque mense Febr. toto nisi perfectis aut reiectis legationibus.

\ciceroSection{2} haec tamen opinio. est populi Romani, a tuis invidis atque obtrectatoribus nomen inductum fictae religionis, non tam ut te impediret, quam ut ne quis propter exercitus cupiditatem Alexandriam vellet ire. dignitatis autem tuae nemo est quin existimet habitam esse rationem ab senatu; nemo est enim, qui nesciat, quo minus discessio fieret, per adversarios tuos esse factum; qui nunc populi nomine, re autem vera sceleratissimo tribunorum latrocinio si quae conabuntur agere, satis mi provisum est, ut ne quid salvis auspiciis aut legibus aut etiam sine vi agere possent.

\ciceroSection{3} ego neque de meo studio neque de non nullorum iniuria scribendum mihi esse arbitror; quid enim aut me ostentem, qui, si vitam pro tua dignitate profundam, nullam partem videar meritorum tuorum adsecutus, aut de aliorum iniuriis querar, quod sine summo dolore facere non possum? ego tibi a vi, hac praesertim imbecillitate magistratuum, praestare nihil possum; vi excepta possum confirmare te et senatus et populi Romani summo studio amplitudinem tuam retenturum.
